\chapter{Los visitantes}
\label{ch:tourism}

\section{La industria que Cuba construyó, y lo que costó}

Cuba reconstruyó el turismo casi desde cero después de 1990 ---340.000 visitantes ese
año, 4,7 millones en 2018--- y por cualquier medida convencional eso es un éxito. El
problema es lo que pasó después, y es el ejemplo más nítido de este libro de inversión
asignada por poder antes que por rentabilidad.

A lo largo de la década hasta 2024 el Estado cubano, principalmente a través del
conglomerado militar GAESA y su brazo hotelero, construyó habitaciones de hotel a un
ritmo sostenido y creciente mientras la ocupación se derrumbaba.

Vale la pena exponer las cifras de 2024 completas, y en proporciones antes que en pesos,
porque las proporciones son inmunes al problema del tipo de cambio de la nota sobre las
fuentes. De la inversión nacional total de ese año, los servicios empresariales, las
actividades inmobiliarias y de alquiler se llevaron 24.900 millones de pesos y los hoteles
y restaurantes otros 11.900 millones ---juntos, el 37,4 por ciento de todo lo que el país
invirtió---. La agricultura, la ganadería y la silvicultura se llevaron 2.600 millones, o
el 2,7 por ciento. La salud pública y la asistencia social se llevaron 2.000 millones,
menos de una quinta parte de lo que fue a hoteles y restaurantes solamente. Como lo
formuló el economista Pedro Monreal, la inversión agrícola fue catorce veces menor que la
del turismo \citep{oneiinv}.\unv{}

Esa es la asignación de un país que importa la mayor parte de su comida, es incapaz de
iluminar sus ciudades, pide leche para niños al Programa Mundial de Alimentos, y opera sus
hoteles existentes a tasas de ocupación reportadas en torno a una cuarta parte. Se
abrieron hoteles nuevos hacia un mercado sin huéspedes mientras las estaciones de bombeo
que abastecen a esos mismos hoteles fallaban por falta de piezas.

Ninguna señal de mercado produjo eso, y ninguna señal de mercado podría haberlo hecho. Lo
produjo una institución capaz de asignarse capital a sí misma, sin publicar un
rendimiento, sin competir por los fondos, y sin que nadie estuviera en posición de decir
que no. El capítulo~\ref{ch:integrity} lo usa como ejemplo desarrollado, porque demuestra
que el problema de inversión de Cuba no es en primer lugar una escasez de capital. Es que
el capital que existe lo asigna quien lo controla, hacia lo que controla.

Otros dos rasgos de la industria tal como se construyó importan para lo que sigue.

Fue un \emph{enclave}. La forma dominante es el resort de playa todo incluido en Varadero
y los cayos del norte, físicamente separado del país que lo rodea, abastecido en gran
medida con comida importada, atendido por trabajadores pagados en pesos a través de una
agencia estatal que se queda con el diferencial cambiario y ---hasta 2008--- legalmente
cerrado a los cubanos. La evidencia de todo el Caribe sobre este modelo es que la
proporción del gasto del visitante que se queda en el país anfitrión es baja, citada
habitualmente entre una quinta y una tercera parte para los todo incluido, con el resto
acumulándose para el operador turístico extranjero, la aerolínea extranjera y la cadena de
suministro importada.\unv{} Cuba adoptó el modelo que más se filtra.

Y arrastra una memoria. El capítulo~\ref{ch:inheritance} describió La Habana de los
años cincuenta: casinos, crimen organizado, un distrito de placer para extranjeros en un
país cuya propia gente quedaba excluida. La reivindicación fundacional de la revolución
incluía abolir aquello. Toda propuesta de convertir el turismo en industria ancla tiene
que responder a la objeción de que lo recrea, y la respuesta no puede ser que los tiempos
cambiaron. Tiene que ser un diseño en el que los cubanos sean dueños de los negocios, no
estén excluidos de los espacios, y capturen el ingreso.

\section{La estrategia: valor por visitante, no visitantes}

El libro propone una inversión explícita de la métrica con la que se ha manejado la
industria.

Cuba no debe perseguir el número de llegadas. Debe perseguir la \emph{divisa neta
retenida por visitante}, y debe estar dispuesta a aceptar menos visitantes para elevarla.

El razonamiento es en parte ambiental y en parte aritmético. El capítulo~\ref{ch:nature}
expone lo que el turismo de masas le hace a los arrecifes, los humedales y el tejido
colonial, y la ventaja comparativa de Cuba es precisamente aquello que el volumen destruye.
Pero la aritmética basta por sí sola: un visitante alojado en una casa particular, que come
en restaurantes privados y compra a guías privados, deja en manos cubanas varias veces más
que un visitante en un todo incluido, incluso cuando el segundo pagó más por sus
vacaciones. En una isla limitada por la mano de obra y con una red eléctrica que falla,
añadir visitantes es caro ---cada uno necesita agua, electricidad, comida y personal, todo
escaso--- mientras que añadir valor por visitante es casi gratis.

Esto resuelve además la tensión con la restricción de mano de obra del
capítulo~\ref{ch:premises}. Cuba no tiene gente para un Cancún. Sí tiene gente para algo
mejor pagado.

\begin{design}{La economía del visitante}
\label{d:visitor}
\begin{rules}
\rl{1}{La métrica} El objetivo publicado de la política turística es la divisa neta
retenida por noche de visitante, informada anualmente con sus componentes. No se fija ni
se publica meta alguna de llegadas. Lo que un Estado mide es lo que construye, y Cuba
construyó habitaciones porque contaba llegadas.

\rl{2}{Privado por defecto} El alojamiento, la restauración, el guiado, el transporte y el
comercio al visitante quedan abiertos por derecho a prestadores privados y cooperativos,
según el Diseño~\ref{d:commit}.2, con el papel del Estado limitado a normas, seguridad e
inspección. La casa particular y el paladar son la forma por defecto de la industria
turística cubana, no su margen tolerado.

\rl{3}{Empleo y pago directos} Se suprime el arreglo de la agencia empleadora estatal para
este sector como para todos los demás (capítulo~\ref{ch:capital}): los hoteles contratan a
su propio personal y le pagan directamente. El diferencial cambiario que se le quita al
trabajador turístico es la injusticia fundacional de la industria y es además la razón de
que el servicio sea indiferente.

\rl{4}{Compra local} Los negocios de alojamiento por encima de un umbral de tamaño tienen
que publicar la proporción de alimentos comprada en el país, y el gravamen al alojamiento
del Diseño~\ref{d:tax}.7 se fija en dos tasas, aplicándose la menor por encima de un umbral
publicado de compra nacional. Es el instrumento que conecta el capítulo~\ref{ch:land} con
este: el cliente potencial más fiable de la agricultura cubana está parado en un hotel a
dos horas, comprando tomates importados.

\rl{5}{Ningún hotel estatal nuevo} Moratoria a la construcción hotelera estatal hasta que
la ocupación del sistema supere un umbral publicado durante dos años consecutivos. El
capital liberado se redirige según la tabla de financiamiento del
capítulo~\ref{ch:sequence}. Es la mayor reasignación de inversión disponible para Cuba y no
exige ni un dólar de dinero extranjero.

\rl{6}{Abierto a los cubanos} Ninguna instalación, playa, cayo o distrito queda cerrado a
los residentes cubanos, ni en la ley ni en la práctica, y la prohibición es justiciable
según el Diseño~\ref{d:courts}.4.
\end{rules}
\end{design}

\section{Lo que Cuba tiene de verdad}

El activo no es la playa. Todos los países del Caribe tienen playa, varios la tienen mejor,
y competir en eso significa competir en precio con lugares que tienen mejores aeropuertos y
electricidad que funciona.

Lo que Cuba tiene y la competencia no es una combinación que no puede construirse: nueve
sitios del Patrimonio Mundial de la UNESCO; cuatro ciudades coloniales intactas; una
cultura musical que es un destino en sí misma; algunos de los corales mejor conservados del
hemisferio en los Jardines de la Reina; el mayor humedal del Caribe insular en Zapata, con
la lista de aves endémicas que lo acompaña; cinco mil setecientos kilómetros de costa,
buena parte sin desarrollar; y ---cosa inusual e importante--- una tasa de crimen violento
muy baja, que es una ventaja decisiva en una región donde varios competidores tienen la
reputación contraria.

Los productos que se siguen de esa lista son el buceo, la observación de aves y el turismo
de naturaleza, el turismo cultural y musical, el turismo arquitectónico y patrimonial, el
turismo de salud del capítulo~\ref{ch:life}, y los visitantes de larga estancia del
capítulo~\ref{ch:talent} ---todos ellos de mayor valor, menor volumen, menos estacionales y
menos destructivos ambientalmente que los resorts de playa, y todos ellos hoy
subdesarrollados o inexistentes---.

\subsection{El precedente de La Habana Vieja}

Cuba ya construyó la mejor institución patrimonial del Caribe y después la desmanteló, y
el episodio es instructivo en las dos direcciones.

La Oficina del Historiador de la Ciudad de La Habana, bajo Eusebio Leal, recibió en 1993
algo que ningún otro organismo cubano tenía: el derecho a operar empresas comerciales en
La Habana Vieja ---hoteles, restaurantes, tiendas--- y a \emph{retener y reinvertir sus
propios ingresos} en restaurar el distrito. Funcionó. La restauración de La Habana Vieja es
el proyecto físico más exitoso del período cubano, se autofinanciaba, empleó y formó a
residentes locales, restauró vivienda junto a monumentos en vez de desplazar gente, y
produjo un distrito que es hoy el activo turístico individual más fuerte del país.

En 2016 su brazo comercial se transfirió al conglomerado militar, y con él la base de
ingresos independiente de la Oficina.

Las dos mitades pertenecen al diseño. El mecanismo ---una autoridad patrimonial que
conserva sus propios ingresos y responde por el tejido--- es correcto y debería extenderse a
Trinidad, Cienfuegos, Camagüey y Santiago. Y la reversión es la razón por la que hay que
anclarlo según el Diseño~\ref{d:commit} en vez de concederlo, porque se concedió, funcionó
veintitrés años, y se lo quitaron.

\begin{design}{Autoridades patrimoniales}
\label{d:heritage}
\begin{rules}
\rl{1}{Forma} Una autoridad patrimonial estatutaria para cada ciudad histórica, con
jurisdicción sobre la licencia urbanística, las normas de edificación y el espacio público
de su distrito.

\rl{2}{Ingresos} Retiene una proporción definida del gravamen al alojamiento y de los
ingresos por concesión recaudados en su distrito, por ley y no por concesión graciosa, y esa
proporción solo puede alterarse mediante ley.

\rl{3}{Obligación} Un programa de restauración publicado y auditado, y un requisito
vinculante de que una proporción fija de la superficie restaurada siga siendo residencial y
de que los residentes existentes tengan derecho a realojo dentro del distrito. El turismo
patrimonial que vacía un distrito de sus habitantes produce un decorado y pierde aquello por
lo que vinieron los visitantes.

\rl{4}{Independencia} Su brazo comercial no puede transferirse a ninguna otra entidad sin
una ley. Escrito porque es exactamente lo que ocurrió.
\end{rules}
\end{design}

\section{Capacidad de carga}

El capítulo~\ref{ch:nature} desarrolla el argumento ambiental. Lo que corresponde aquí es
que es un argumento económico y no sentimental.

Los arrecifes cubanos, el humedal de Zapata y los cayos son el acervo físico desde el cual se
venden los productos de alto valor de arriba. Un arrecife blanqueado, dinamitado por anclas o
enterrado en sedimento por la construcción costera deja de producir ingresos de forma
permanente, y los ingresos que producía eran los más altos por visitante de la industria. El
buceo en los Jardines de la Reina funciona precisamente porque se excluyó la pesca y los
tiburones siguen ahí; vale más por día de visitante que cualquier resort de playa del país y
existe solo porque alguien trazó una línea en un mapa y la hizo cumplir.

\begin{design}{Límites}
\label{d:limits}
\begin{rules}
\rl{1}{Límites por sitio} Límites diarios y anuales publicados de visitantes para arrecifes,
cayos y áreas protegidas, fijados por el regulador ambiental del capítulo~\ref{ch:nature} y
no por el ministerio de turismo, con el acceso asignado por subasta o sorteo publicados y no
por permiso administrativo.

\rl{2}{Retiro costero} Un retiro de construcción legal respecto de la línea de costa,
aplicado sin excepción, incluidos los operadores estatales y los vinculados a los militares.
El capítulo~\ref{ch:nature} fija la distancia; esta cláusula trata sobre la ausencia de
excepciones, que es donde ha fracasado toda regla de retiro costero del Caribe.

\rl{3}{Ningún pedraplén nuevo} Los pedraplenes que conectan los cayos del norte con tierra
firme han dañado de forma medible la circulación del agua y los ecosistemas que quedan
detrás. No se construyen más, y se evalúa la remediación de los existentes.

\rl{4}{Ponerle precio a la escasez} Donde un sitio esté en su límite, el precio del acceso
sube. Es el mecanismo que convierte una restricción ambiental en ingreso bajo el
Diseño~\ref{d:visitor}.1 en vez de en una cola.
\end{rules}
\end{design}

\section{La enfermedad holandesa, y el problema salarial}

Un sector turístico exitoso que gane divisas, bajo el régimen dolarizado del
capítulo~\ref{ch:money}, elevará los precios y los salarios en los sectores que lo atienden y
estrujará todo lo que compite con importaciones ---que es la agricultura, que es el sector
que el capítulo~\ref{ch:land} necesita hacer crecer---. Es la patología clásica del auge de
recursos y aquí se agrava por la restricción de mano de obra: no hay reserva de trabajadores
para expandirse hacia el sector en auge, de modo que este los saca de los demás.

Cuba ya corrió este experimento. En los años noventa el salario turístico llevó a los médicos
a manejar taxis y a los cirujanos a alquilar cuartos, que es de donde salió la pirámide
invertida del capítulo~\ref{ch:dollarreform}.

\begin{design}{No repetir los años noventa}
\label{d:dutch}
\begin{rules}
\rl{1}{La paga pública es el control} El Diseño~\ref{d:pay} obliga aquí. La paga profesional
pública se referencia contra el salario turístico y privado y se mueve con él. Una estrategia
turística adoptada sin el acuerdo correspondiente de paga pública es una estrategia para
vaciar los hospitales, y este libro no va a proponer una.

\rl{2}{Los ingresos al fondo} El gravamen al alojamiento y los ingresos por concesión son
renta según el Diseño~\ref{d:rentrule} y van al fondo de estabilización e inversión, no al
presupuesto corriente ---lo que a la vez amortigua el impulso de demanda doméstica y evita
que un auge turístico financie compromisos corrientes que una caída no podrá sostener---.

\rl{3}{Encadenamiento antes que expansión} El esfuerzo de política pública va a elevar el
contenido nacional del gasto de cada visitante ---comida, muebles, textiles, artesanía,
guiado--- antes que a elevar el número de visitantes. El encadenamiento es el instrumento que
convierte un sector en auge en uno diversificado, y es la razón de que el
Diseño~\ref{d:visitor}.4 esté escrito como un diferencial fiscal y no como una exhortación.

\rl{4}{Ponderación regional} Los incentivos al desarrollo turístico se ponderan hacia las
provincias orientales, según el Diseño~\ref{d:municipal}.2. Dejada a su aire, la industria se
concentra en La Habana, Varadero y los cayos del norte, y entonces la fórmula de igualación
tiene que hacer todo el trabajo.
\end{rules}
\end{design}
